\newpage
\section{Parametry Ollamy}
\label{sec:parametry}

Każde zapytanie do Ollamy ma obiekt \kod{options}. Wartości używane w aplikacji są w \kod{config.yaml}: sekcja
\kod{generation} dla moderacji (krótkie, powtarzalne odpowiedzi JSON) i sekcja \kod{chat} dla rozmowy (dłuższy,
trochę swobodniejszy tekst).

\begin{table}[H]
\centering
\small
\begin{tabularx}{\linewidth}{@{}l l X@{}}
\toprule
Parametr & Wartość & Znaczenie \\
\midrule
\kod{temperature} & 0,1 (czat 0,3) & dzieli logity przed losowaniem; mała wartość -- prawie zawsze wygrywa najbardziej
  prawdopodobny token, duża -- rozkład się spłaszcza i tekst jest bardziej różnorodny \\
\kod{top\_p} & 0,9 & \emph{nucleus sampling}: losowanie tylko z najmniejszego zbioru tokenów, których
  prawdopodobieństwa sumują się do \kod{top\_p} \\
\kod{top\_k} & 40 & losowanie tylko z \kod{top\_k} najbardziej prawdopodobnych tokenów; \kod{top\_k = 1} to
  dekodowanie zachłanne (\emph{greedy}) z laboratorium 1 \\
\kod{seed} & 42 & stałe ziarno generatora losowego -- ten sam prompt i te same opcje dają tę samą odpowiedź \\
\kod{num\_ctx} & 4096 (czat 8192) & okno kontekstu w tokenach; muszą się w nim zmieścić prompt, przykłady,
  paragrafy i odpowiedź \\
\kod{num\_predict} & 300 (czat 700) & maksymalna liczba wygenerowanych tokenów \\
\kod{repeat\_penalty} & 1,1 & obniża prawdopodobieństwo tokenów, które już wystąpiły (przeciw pętlom, które
  robił GPT-2) \\
\kod{stop} & brak & napisy, na których generowanie się kończy \\
\kod{format} & schemat JSON & ogranicza generowanie do poprawnego JSON-a zgodnego ze schematem \\
\kod{keep\_alive} & 30 min & jak długo model zostaje w pamięci po ostatnim zapytaniu \\
\bottomrule
\end{tabularx}
\caption{Parametry generowania w Ollamie.}
\end{table}

\subsection{Eksperyment}

Skrypt \kod{scripts/parameter\_experiment.py} wysyła do \kod{qwen2.5:7b} te same zapytania z różnymi opcjami i zapisuje
wyniki w \kod{examples/parameters/parameters.json}. Ziarno jest w nim domyślnie wyłączone, żeby było widać naturalną
zmienność modelu. Wyniki jednego uruchomienia:

\begin{table}[H]
\centering
\small
\begin{tabularx}{\linewidth}{@{}>{\raggedright\arraybackslash}p{0.4\linewidth} X@{}}
\toprule
Test & Wynik \\
\midrule
\kod{temperature} 0 / 0,7 / 1,5, ocena komentarza ,,Kurde, znowu zapomniałem średnika'' 5 razy (schemat JSON)
  & za każdym razem \kod{review}, stopień 1 -- przy wymuszonym schemacie i jasnym prompcie decyzja nie zmieniła się
    nawet przy $T = 1{,}5$ \\
\kod{seed} przy $T = 0{,}9$, 3 odpowiedzi tekstowe & z ziarnem 42: 1 różna odpowiedź na 3, bez ziarna: 3 różne na 3 \\
\kod{top\_p} 0,1 / 1,0 przy $T = 1{,}0$, 5 odpowiedzi & 1 / 5 różnych odpowiedzi \\
\kod{top\_k} 1 / 100 przy $T = 1{,}0$, 5 odpowiedzi & 1 / 5 różnych odpowiedzi \\
\kod{num\_predict = 15} & odpowiedź urwana w połowie: \kod{\{ "toxic": false, "categories": ["none"], "severity} --
  takiego JSON-a nie da się odczytać \\
\kod{format}: brak / \kod{"json"} / schemat, 5 komentarzy & poprawne oceny: 0 / 0 / 5 na 5; bez schematu model
  zwracał skrócony JSON, np. \kod{\{"toxic": true, "reason": "insult"\}} \\
\kod{num\_ctx} 2048 / 16384, krótki prompt & 1,26 s / 1,26 s -- przy krótkim prompcie rozmiar okna nie wpływa na czas,
  ma znaczenie dopiero przy długim wejściu (czat o pliku) i dla zajętości pamięci \\
\kod{stop = [","]} & odpowiedź kończy się przed pierwszym przecinkiem \\
\bottomrule
\end{tabularx}
\caption{Wpływ parametrów -- wyniki \kod{scripts/parameter\_experiment.py}.}
\end{table}

Najważniejszy wniosek jest taki, że o powtarzalności decyzji bardziej decyduje wymuszony schemat (\kod{format}) niż
sama temperatura. Temperatura, \kod{top\_p} i \kod{top\_k} wyraźnie zmieniają za to swobodny tekst -- dlatego w czacie
temperatura jest trochę wyższa niż przy moderacji, a ziarno jest stałe w obu przypadkach.
